\week{W1}{Mechanisms, DoF, Grübler–Kutzbach}{1}
\begin{multicols}{3}
\hd{Words and chain types \pg{18--26}}
\textbf{Kinematics} = geometry of motion, no forces. \textbf{Dynamics} = statics (no acceleration) + kinetics (acceleration). \textbf{DoF} = number of independent parameters fixing the configuration \pg{27}. \textbf{Mobility} $M$ = DoF left after all joints act = minimum number of actuators \pg{39}.\\
\textbf{Closed chain}: has a loop. \textbf{Open}: none. \textbf{Serial}: open, every link but first and last touches two others. \textbf{Parallel} \added: end effector joined to base by several chains (closed loops): delta, Stewart.\\
Why loops \pg{26}: motors on the base so moving mass is small; larger load (legs share); interesting paths. A 6-joint arm is open and serial; a four-bar is closed and not serial.

\hd{Rigid-body DoF and joints \pg{27--36}}
Rigid body: \textbf{3 DoF in 2D}, \textbf{6 DoF in 3D}. Point in 2D: 2. Pendulum in 2D: 1.\par
\begin{tabularx}{\linewidth}{@{}Xccc@{}}\toprule
Joint & $F$ & cuts 3D & cuts 2D\\\midrule
Revolute R, Prismatic P & 1 & 5 & 2\\
Helical H & 1 & 5 & --\\
Cylindrical C, Universal U & 2 & 4 & --\\
Spherical S & 3 & 3 & --\\\bottomrule
\end{tabularx}
“Cuts” $=D-F$ constraints removed. Planar mechanisms use only R and P. Two free planar links $=3+3=6$; one R pin removes $3-1=2$, leaving 4 \pg{36} \chk.
\begin{trap}
Joint DoF $\neq$ mechanism mobility. Helical $=1$ (pitch ties turn to slide). A 2D body is 3, not 2.
\end{trap}

\hd{Grübler–Kutzbach \pg{39--41}}
\begin{fbx}[Mobility, form 2 (use this one)]
\[M=D\,(N-J-1)+\sum_{i=1}^{J}F_i\]
$D=3$ planar, $6$ spatial; $N$ = links \emph{including ground}; $J$ = joints; $F_i$ = DoF of joint $i$.
\use{any count. Serial chain: $N=J+1$ so $M=\sum F_i$.}
\end{fbx}
\begin{fbx}[Form 1 \pg{40} (same sum, grouped)]
\[M_{3D}=6(N{-}1)-5P_1-4P_2-3P_3-2P_4-P_5,\quad M_{2D}=3(N{-}1)-2P_1-P_2\]
$P_i$ = number of joints with $i$ DoF.
\use{cross-check; each joint removes $D-F$.}
\end{fbx}
\begin{recipe}[Mobility count (W13 Q6, Q7)]
\label{w1:gk}
\begin{enumerate}
\item $D$: all axes $\perp$ page $\Rightarrow3$. Any S, U, C or “still 3D” $\Rightarrow6$.
\item Number links, \textbf{ground = 0}. Split every prismatic joint into two links (cylinder, rod).
\item List every joint with its $F$; get $J$ and $\sum F_i$.
\item Write $M=D(N-J-1)+\sum F_i$, then substitute.
\item Check: $M\ge1$ for a machine; if $M\le0$ ask whether constraints are dependent.
\end{enumerate}
\end{recipe}

\columnbreak
\sh{Counts to know}
\begin{tabularx}{\linewidth}{@{}Xcccc@{}}\toprule
Mechanism & $D$ & $N$ & $J$ & $M$\\\midrule
Planar serial RPR \pg{42} & 3 & 4 & 3 & 3\\
Four-bar \pg{44} (nb Q2) & 3 & 4 & 4 & 1\\
3-RPR parallel \pg{46} & 3 & 8 & 9 & 3\\
Five-bar \added & 3 & 5 & 5 & 2\\
Slider-crank \added & 3 & 4 & 4 & 1\\
Stewart 6-UPS \pg{48} & 6 & 14 & 18 & 6\\
Delta, plain & 6 & 11 & 15 & 9 $\to$ 3\\\bottomrule
\end{tabularx}
\begin{worked}[four-bar, 3-RPR]
Four-bar: $3(4-4-1)+4=1$; form 1: $3(4-1)-2\cdot4=1$. 3-RPR: $N=1+1+3\cdot2=8$, $3(8-9-1)+9=3$. \chk
\end{worked}
\begin{worked}[Stewart, D = 6]
$N=2+6\cdot2=14$, $J=18$, $\sum F=6(2+1+3)=36$: $6(14-18-1)+36=6$. One leg alone: $6(4-3-1)+6=6$. Six SPS legs give 12 (leg spins): the U blocks them \added. \chk
\end{worked}
\begin{worked}[delta, W13 Q7]
$N=2+3\cdot3=11$, $J=3(1+4)=15$, $\sum F=3(1+4\cdot3)=39$. Plain: $6(11-15-1)+39=\mathbf 9$. Six rods spin freely: $R_M=6$, $R_C=0$:
\[M=D(N{-}J{-}1)+\sum F_i+R_C-R_M=9+0-6=\mathbf 3\]
Plain 9 if only plain G--K is quoted; 3 if the extended equation is given or the real platform DoF is asked. \chk
\end{worked}
\begin{trap}
(1) Forgetting ground: 3-RPR gives $0$, Stewart $0$. (2) $D=6$ on a planar drawing: four-bar $-2$. (3) $D=3$ for Stewart: $21$. (4) Constraints removed $=F$ (not $D-F$): four-bar $5$. (5) Prismatic counted as one link. (6) Stopping at 9 for the delta. \chk
\end{trap}

\hd{Where G--K fails \pg{49--51}}
Three equal parallel legs (R at both ends), ground--legs--body: $N=5$, $J=6$, $M=3(5-1)-2\cdot6=0$, \textbf{but it moves: 1 DoF}. Rank of constraint Jacobian is 5, not 6: the third leg repeats a constraint. Delete a middle leg: $N=4$, $J=4$, $M=1$ (four-bar). \chk
\begin{recipe}[“Did you notice anything odd?”]
\label{w1:odd}
\begin{enumerate}
\item Do the count, show it ($M=0$).
\item State: the mechanism visibly moves, 1 DoF.
\item Reason (slide wording): G--K assumes \emph{independent} constraints; here one leg is redundant.
\item Fix: remove one leg, recount ($N=4,J=4,M=1$).
\item Caveat: only for equal, parallel legs; with unequal legs it locks and $0$ is right.
\end{enumerate}
\end{recipe}
Hoberman \pg{51}: 24 links, 36 R: $3(24-36-1)+36=-3$ ($D{=}3$), $-42$ ($D{=}6$); both negative, yet it moves \added. \chk
\begin{trap}
False: “$M=0$ means it cannot move”; “G--K always gives the true mobility”. G--K is exact only for independent constraints; it can under-count (p.49) or over-count (delta, local spins).
\end{trap}

\columnbreak
\hd{Actuation \pg{37--38}}
Controllable DoF = actuated. Compare actuators with total DoF:
\begin{tabularx}{\linewidth}{@{}lX@{}}\toprule
equal & \textbf{fully actuated}\\
fewer & \textbf{under-actuated}: cheaper, simpler, lighter, but not every DoF can be commanded \pg{38}\\
more & \textbf{redundant} (slide wording)\\\bottomrule
\end{tabularx}
6R arm, 6 motors: fully; 4 motors: under. Stewart, 6 actuators: fully. Delta, 3 motors, real $M=3$: fully (plain 9 would wrongly say under). \chk
\begin{slideerr}[redundant]
Slide 37 says “controllable $>$ total $\Rightarrow$ redundant”, yet nb cell 8 calls a 7-DoF arm redundant (7 motors = 7 DoF). Standard meaning \added: more joint DoF than the task needs (7 $>$ 6 for a 3D pose; planar 3R positioning a point: $3>2$). Write “redundant: more joint DoF than the 6-DoF task” and say which meaning you use.
\end{slideerr}

\hd{Workspace \pg{52--57}}
\textbf{Reachable} $W_R$: reached in $\ge1$ orientation. \textbf{Dexterous} $W_D$: reached in \emph{all} orientations. Set by link geometry and joint limits \pg{52}.
\begin{fbx}[Relation \pg{57}]
\[W_{D}\subseteq W_{R}\qquad(\text{i.e. }W_R\supseteq W_D)\]
\use{answer “$W_R\supseteq W_D$”. $\subset$ strict and $\supset$ are traps (can be equal, e.g.\ both empty). $W_R\subseteq W_D$ is backwards.}
\end{fbx}
\begin{fbx}[Planar 2R and 3R reach \added]
2R: $|L_1-L_2|\le r\le L_1+L_2$. $L=(0.5,0.3)$: annulus $0.2\le r\le0.8$\,m.\\
3R, equal links $L$: $W_R$ disc $r\le3L$; $W_D$ disc $r\le L$ (circle of radius $L$ about the tip must fit in $0\le|\mathbf w|\le2L$). $L=0.3$: $0.9$\,m and $0.3$\,m. \chk
\use{fewer DoF than the pose (3 planar, 6 spatial) $\Rightarrow$ $W_D$ empty (any 2R arm).}
\end{fbx}
\begin{tabularx}{\linewidth}{@{}llX@{}}\toprule
Arm & Joints & Workspace shape \pg{53--56}\\\midrule
PUMA & RRR & thick spherical shell, wedge cut by joint limits\\
SCARA & RRP & flat ring (annulus), thickened by vertical stroke\\
Cylindrical & RPP & thick-walled cylinder sector\\
Cartesian & PPP & box, edges = the three strokes\\\bottomrule
\end{tabularx}
Shapes read from the slide figures.

\hd{What the unit covers \pg{58--62}}
FK (W2): $\mathbf q\to$ pose. IK (W3): pose $\to\mathbf q$. Velocity (W6), dynamics (W7--W9), trajectories (W5). Pose = 3 + 3 numbers = a $4\times4$ homogeneous matrix. Frames: {B} base, {W} wrist, {T} tool, {S} station, {G} goal, {C} camera \pg{61}.

\hd{Practical (optional)}
\texttt{test.py}: CoppeliaSim remote API, drives joint \texttt{/Joint1} to $60^\circ$ for 3\,s; one revolute joint. No saved outputs.
\end{multicols}
